\section{Introduction}
\label{sec:intro}

A central question in the economics of artificial intelligence is whether autonomous pricing
agents---deployed independently by competing firms, with no shared reward function and no
explicit coordination protocol---converge to supra-competitive (collusive) outcomes purely
through repeated interaction in a market. \citet{calvano_etal_2020_ai_algorithmic_pricing_collusion} demonstrated that two
independent $Q$-learners in a repeated Bertrand duopoly converge to prices well above the
static Nash equilibrium and approach the monopoly benchmark, with terminal prices distributed
across multiple distinct attractors as a function of initial $Q$-matrix and exploration
parameters. Subsequent empirical work \citep{assad_etal_2024_algorithmic_pricing_competition_empirical,fish2024algorithmicCollusion}
has extended these findings to richer agent classes and field settings, and a parallel
regulatory literature \citep{ezrachi2016virtual,asker2021ai_pricing} has framed the
phenomenon as a first-order antitrust concern: existing competition law typically requires
evidence of explicit agreement, which self-learning agents do not produce.

The arrival of large language model (LLM) agents introduces a third algorithmic family with
qualitatively different properties from $Q$-learners. LLMs carry no shared reward signal, no
explicit value function, no online weight updates, and no exploration parameter; their
``learning'' within a pricing game is in-context, mediated through textual reasoning and
prompt history. Whether such agents reproduce the basin-of-attraction structure that
characterizes $Q$-learner outcomes
\citep{waltman_kaymak_2008_qlearning_cournot,wunder_littman_babes_2010_qlearning_basin}
is largely uncharacterized in the literature.\footnote{Concurrent work
\citep{Tailor_2025_whisper} introduces a multi-domain collusion-detection benchmark spanning
several economic and adversarial settings; our scope is the basin-of-attraction
characterization at fixed Bertrand-pricing setting and is complementary in scope.}

In this paper we report Stage 2b of a pre-registered pilot study of LLM-agent Bertrand
pricing dynamics. Stage 2b extends the horizon of an earlier 25-round Stage 2 protocol to
50 rounds across two conditions: GATE-5 ($n=5$ agents, 30 repetitions) and GATE-2 ($n=2$
agents, 15 repetitions). The pre-registered analysis pipeline is locked on OSF Registries
prior to data collection and timestamped with RFC 3161 trusted timestamps; the same pipeline
is independently re-implemented in Python with SciPy reference statistical primitives, and
the full pipeline runs against the locked dataset are mutually verified to within numerical
tolerance.

Our headline finding is an asymmetry between two conditions whose static profit measure is
nearly identical but whose dynamic convergence properties differ sharply: $\Delta_{\text{profit}}$
at convergence is $0.4932$ ($n=5$) vs.\ $0.4522$ ($n=2$)---both capturing roughly half the
Nash-to-Monopoly gap---while convergence within the 50-round horizon occurs in $30/30$
repetitions for $n=5$ and only $1/15$ repetitions for $n=2$. The $n=2$ non-convergent
population exhibits sustained regime drift that we characterize, in a pre-committed secondary
analysis, as basin-stable but not statically converged. We argue this asymmetry is consistent
with a wider basin of attraction at smaller agent counts: a larger fraction of trajectories
land in the supra-competitive regime under $n=5$, while $n=2$ trajectories occupy a wider
distribution of within-basin behaviors that have not equilibrated by round 50. The mechanism
behind this asymmetry is not the cross-sectional-signal reading we initially proposed: a
pre-registered text-based test (\S\ref{sec:discussion-asymmetry-cot}) reversed that reading
on Stage 2b data; a refined reading --- per-peer attention density at $n = 2$ exceeds the
rate at $n = 5$ even after controlling for total CoT word count --- is supported by
Stage 2b data and locked for Stage 4 cross-model generalization.

As a co-equal headline, the bimodal $k = 2$ basin structure that \citet{Cao_2026_llm_collusion}
reports for LLM-agent pricing is recovered at $n = 5$ ($7/10$ gap-statistic resamples;
$30/30$ basin-stable) and not recovered at $n = 2$ ($k = 1$ in $10/10$ resamples)
within the Stage 2b $n_{\text{rep}} = 15$ sample-size envelope --- the basin signature
is conditional on agent count, not a fixed property of LLM-agent pricing.

Our contributions are:

\begin{enumerate}
    \item \textbf{Empirical:} A pre-registered, multi-rep characterization of LLM-agent
        Bertrand-pricing dynamics under two agent counts, with full per-repetition
        reporting (\S\ref{sec:results}). To our knowledge, this is the first such study
        with a pre-committed RFC 3161-timestamped analysis pipeline.

    \item \textbf{Methodological:} A verification protocol consisting of (a) cryptographic
        pre-registration as the foundation of analytic-claim integrity, (b) procedural
        cross-toolchain reproduction (TypeScript primary; Python/SciPy verifier) as a
        reinforcing layer, and (c) adversarial multi-persona expert-panel review as a
        further reinforcing layer. The protocol is hierarchical, not symmetric: the
        cryptographic lock is the strongest integrity primitive available for analysis-plan
        locking; the verification and adversarial layers reinforce it but cannot substitute
        for it (\S\ref{sec:methods}).

    \item \textbf{Conceptual:} A pre-registered multi-condition cross-repetition
        characterization showing that LLM bimodal terminal-state structure
        \citep{Cao_2026_llm_collusion} is conditional on agent count---recovered at $n = 5$
        ($k = 2$ in $7/10$ gap-statistic resamples; $30/30$ basin-stable), and not
        recovered at $n = 2$ within the Stage 2b $n_{\text{rep}} = 15$ sample-size envelope
        (\S\ref{sec:limitations-envelope}). The cross-sectional-signal mechanism for this
        conditionality, initially proposed as reading \emph{(a)} of the basin-width
        hypothesis, was empirically reversed by a pre-registered text-based test
        (\S\ref{sec:discussion-asymmetry-cot}), and a refined reading \emph{(a$'$)} ---
        per-peer attention density at $n = 2$ exceeds the rate at $n = 5$ even controlling
        for total CoT word count --- is supported by Stage 2b data via the pre-registered
        talk-volume-controlled residualization. The basin-width-by-agent-count hypothesis is
        locked for Stage 3 falsification testing (\S\ref{sec:discussion-falsifiability})
        and refined reading \emph{(a$'$)} for Stage 4 cross-model generalization
        (\S\ref{sec:discussion-asymmetry-cot}).
\end{enumerate}

We frame the Stage 2b outcomes as partial-coordination---supra-competitive prices reached
without explicit agreement---per the regulator-relevant case identified by
\citet{ezrachi2016virtual,asker2021ai_pricing}; \S\ref{sec:discussion-implications}
elaborates with full citations and the welfare-effect-adjudication scope conditions.

The pre-registered Stage 3 design expands the repetition count to $n \geq 30$ for both
conditions, addresses the basin-stability single-rep fragility that the $n=15$ Stage 2b
sample cannot eliminate, and instruments the runner trace schema to isolate LLM-serving
variability \citep{kwon2023pagedattention,zheng2023sglang,holmes2024deepspeedfastgen} from
reasoning-depth variability \citep{snell2024testtimecompute}---a confound that this Stage 2b
report explicitly does not adjudicate (\S\ref{sec:limitations}).

We position this work as a multi-agent LLM-safety contribution; classical industrial-organization
equilibrium-selection refinements (e.g., welfare-loss quantification, full Stigler-comparative-statics
adjudication, Cournot/asymmetric-cost robustness) are deferred to companion work that builds on
the Stage 3 expansion (\S\ref{sec:limitations-bertrand}).
